%% TODO(redação final): redigir em parágrafo único, entre 150 e 500 palavras,
%% somente depois de consolidados resultados, discussão e conclusão.
\begin{resumo}
Esta dissertação investiga como técnicas de avaliação de sistemas de geração aumentada por recuperação podem ser organizadas e aplicadas à análise da recuperação documental, da geração de respostas e da verificabilidade das fontes em um sistema voltado à consulta de discursos da 56\textordfeminine{} Legislatura do Senado Federal. A pesquisa, de natureza aplicada, exploratória e descritiva, adota abordagem multimétodos e dialoga com a \textit{Design Science Research}. O protocolo avaliativo articula métricas de recuperação, rubrica de avaliação das respostas, julgamento humano independente, avaliação automatizada baseada no paradigma \textit{LLM-as-a-judge} e análise qualitativa de erros. O estudo parte de uma prova de conceito previamente implementada e amplia sua bateria de perguntas, seu conjunto documental de referência e seus controles experimentais. % TODO: completar com o corpus e o período efetivamente analisados, os principais resultados e as conclusões finais.
\end{resumo}
